% Generated draft: requires mathematical and source review.
\documentclass[11pt]{article}
\usepackage[margin=25mm]{geometry}
\usepackage{fontspec}
\setmainfont{Latin Modern Roman}
\usepackage{amsmath,amssymb,mathtools,unicode-math}
\setmathfont{Latin Modern Math}
\usepackage{enumitem,needspace,xurl,hyperref}
\setlength{\parindent}{0pt}
\setlength{\parskip}{5pt}
\setlength{\emergencystretch}{3em}
\newcommand{\N}{\mathbb{N}}
\newcommand{\Z}{\mathbb{Z}}
\newcommand{\Q}{\mathbb{Q}}
\newcommand{\R}{\mathbb{R}}
\newcommand{\C}{\mathbb{C}}
\newcommand{\F}{\mathbb{F}}
\newcommand{\A}{\mathbb{A}}
\newcommand{\PP}{\mathbb{P}}
\newcommand{\E}{\mathbb{E}}
\newcommand{\eps}{\varepsilon}
\newcommand{\dd}{\,\mathrm{d}}
\begin{document}
\section*{UnsolvedMath ID: 701078}
\textbf{Source problem code:} catalogue entry
\subsection{Definitions and mathematical statement}
Let \(p\) be a prime and write \(\F_p=\Z/p\Z\). Interpret \(N(\F_p)\) as the Nottingham group
\[
N(\F_p)=\left\{x+\sum_{n=2}^{\infty}b_nx^n:b_n\in\F_p\right\},
\]
with formal composition as its group operation and the \(x\)-adic topology. Thus a neighbourhood basis of the identity consists of the subgroups of series congruent to \(x\) modulo \(x^m\), for \(m\geq2\). For each positive integer \(r\), let
\[
H_r=\left\{x\left(1+\sum_{i=1}^{\infty}a_i x^{p^r i}\right):a_i\in\F_p\right\}\subseteq N(\F_p),
\]
which the supplied source describes as a closed subgroup, equipped with the induced topology.

For a profinite group \(G\), its \(p\)-cohomological dimension \(\mathrm{cd}_p(G)\) is the least integer \(d\geq0\) such that
\[
H^n_{\mathrm{cts}}(G,M)=0
\quad\text{for every }n>d\text{ and every discrete }p\text{-primary torsion }G\text{-module }M
\]
with continuous action; set \(\mathrm{cd}_p(G)=\infty\) if no such \(d\) exists.

\textbf{Problem (I. B. Fesenko, Kourovka Notebook 14.94).} Determine
\[
\mathrm{cd}_p(H_r)
\qquad\text{for each integer }r\geq1.
\]
The source does not explicitly specify or quantify the prime \(p\), nor define \(N(\F_p)\). The formulation above uses the standard Nottingham-group interpretation and treats \(p\) as an arbitrary fixed prime; any intended restriction on \(p\) requires source confirmation. Malformed braces in the supplied coefficient-field notation have been repaired to \(\Z/p\Z\). The supplied registry labels the problem open; that label is not an independently verified mathematical assessment.
\subsection{Short English statement}
For each positive integer \(r\), determine the continuous \(p\)-cohomological dimension of the closed subgroup of the Nottingham group over \(\F_p\) consisting of series \(x(1+\sum_{i\geq1}a_i x^{p^r i})\). The prime range and the standard interpretation of the ambient group require source confirmation.
\subsection{Sources}
\begin{itemize}
\item Supplied local UnsolvedMath registry record: permanent ID \(701078\), exact source problem code \texttt{catalogue entry}, title \emph{Kourovka Notebook Problem 14.94}; upstream revision \texttt{d371ee917ac9647fbb06cd6fc529deec6f3dda82}.
\item \emph{The Kourovka Notebook (Unsolved Problems in Group Theory)}, original problem 14.94, attributed in the supplied record to I. B. Fesenko, proposed in 1999 and associated with the 14th issue (1999).
\item The supplied background identifies the 21st issue edition (2026), \texttt{arXiv:1401.0300v48}. The supplied collection description names E. I. Khukhro and V. D. Mazurov as editors. These bibliographic details are supplied attribution and have not been independently verified.
\end{itemize}
\end{document}
